\begin{thebibliography}{99}
\item[]\textit{Stochastic semigroups and positive realization}
\bibitem{Flor} P. Flor, On groups of non-negative matrices, \emph{Compositio Mathematica} \textbf{21} (1969), no. 4, 376--382.
\bibitem{Schwarz58} \v S. Schwarz, On the structure of the semigroup of stochastic matrices, \emph{Publications of the Mathematical Institute of the Hungarian Academy of Sciences, Series A} \textbf{9} (1964), 297--311.
\bibitem{BF04} L. Benvenuti and L. Farina, A tutorial on the positive realization problem, \emph{IEEE Transactions on Automatic Control} \textbf{49} (2004), no. 5, 651--664.
\bibitem{Heller58} A. Heller, On stochastic processes derived from Markov chains, \emph{Annals of Mathematical Statistics} \textbf{36} (1965), no. 4, 1286--1291.
\item[]\textit{Automata, weighted series, and recognition}
\bibitem{Rabin} M. O. Rabin, Probabilistic automata, \emph{Information and Control} \textbf{6} (1963), no. 3, 230--245.
\bibitem{Paz} A. Paz, \emph{Introduction to Probabilistic Automata}, Academic Press, New York, 1971.
\bibitem{BR} J. Berstel and C. Reutenauer, \emph{Noncommutative Rational Series with Applications}, Encyclopedia of Mathematics and its Applications 137, Cambridge University Press, Cambridge, 2011.
\bibitem{FP} N. Fijalkow and C. Paperman, Monadic second-order logic with arbitrary monadic predicates, arXiv:1709.03117 (2017).
\bibitem{BP} A. Brodsky and N. Pippenger, Characterizations of 1-way quantum finite automata, \emph{SIAM Journal on Computing} \textbf{31} (2002), no. 5, 1456--1478.
\bibitem{AF57} A. Ambainis and R. Freivalds, 1-way quantum finite automata: strengths, weaknesses and generalizations, arXiv:quant-ph/9802062 (1998).
\bibitem{FSZ} Y. Feng, L. Song, and L. Zhang, Distribution-based bisimulation and bisimulation metric in probabilistic automata, arXiv:1512.05027 (2015).
\bibitem{Steinberg} B. Steinberg, Topological dynamics and recognition of languages, arXiv:1306.1468 (2013).
\item[]\textit{Compact groups and orbit geometry}
\bibitem{Folland} G. B. Folland, \emph{A Course in Abstract Harmonic Analysis}, second ed., Chapman \& Hall/CRC, Boca Raton, 2016.
\bibitem{Hall} B. C. Hall, \emph{Lie Groups, Lie Algebras, and Representations: An Elementary Introduction}, second ed., Graduate Texts in Mathematics 222, Springer, Cham, 2015.
\bibitem{Szarek57} S. J. Szarek, Metric entropy of homogeneous spaces, \emph{Banach Center Publications} \textbf{43} (1998), 395--410.
\bibitem{BG2012} J. Bourgain and A. Gamburd, A spectral gap theorem in $SU(d)$, \emph{Journal of the European Mathematical Society} \textbf{14} (2012), no. 5, 1455--1511.
\bibitem{BreuillardGelander57} E. Breuillard and T. Gelander, On dense free subgroups of Lie groups, \emph{Journal of Algebra} \textbf{261} (2003), no. 2, 448--467.
\bibitem{BreuillardGelanderCorrection57} E. Breuillard and T. Gelander, On dense free subgroups of Lie groups---revisited, arXiv:2605.20568 (2026).
\item[]\textit{Algebraic computation and complexity}
\bibitem{BPR} S. Basu, R. Pollack, and M.-F. Roy, \emph{Algorithms in Real Algebraic Geometry}, second ed., Algorithms and Computation in Mathematics 10, Springer, Berlin, 2006.
\bibitem{BJKP} V. D. Blondel, E. Jeandel, P. Koiran, and N. Portier, Decidable and undecidable problems about quantum automata, \emph{SIAM Journal on Computing} \textbf{34} (2005), no. 6, 1464--1473.
\bibitem{DJK} H. Derksen, E. Jeandel, and P. Koiran, Quantum automata and algebraic groups, \emph{Journal of Symbolic Computation} \textbf{39} (2005), nos. 3--4, 357--371.
\bibitem{deGraaf} W. A. de Graaf, Computing the Zariski closure of a finitely generated matrix group, arXiv:2505.02577v2 (2026).
\bibitem{NPSW} K. Nosan, A. Pouly, S. Schmitz, M. Shirmohammadi, and J. Worrell, On the computation of the Zariski closure of finitely generated groups of matrices, arXiv:2106.01853v6 (2025).
\bibitem{Shitov} Y. Shitov, A universality theorem for nonnegative matrix factorizations, arXiv:1606.09068v2 (2018).
\bibitem{Ramsey} F. P. Ramsey, On a problem of formal logic, \emph{Proceedings of the London Mathematical Society} (2) \textbf{30} (1930), no. 1, 264--286.
\bibitem{SJR2004}
S.~Singh, M.~R.~James, and M.~Rudary,
\emph{Predictive state representations: A new theory for modeling dynamical systems},
Proceedings of the Twentieth Conference on Uncertainty in Artificial Intelligence,
AUAI Press, 2004, 512--519. See Theorems~2 and~6; arXiv:1207.4167.
\item[]\textit{Entropy transport and explicit spherical spectral bounds}
\bibitem{PW61} Y. Polyanskiy and Y. Wu, Wasserstein continuity of entropy and outer bounds for interference channels, \emph{IEEE Transactions on Information Theory} \textbf{62} (2016), no. 7, 3992--4002. DOI: 10.1109/TIT.2016.2562630.
\bibitem{LPS61} A. Lubotzky, R. Phillips, and P. Sarnak, Hecke operators and distributing points on the sphere I, \emph{Communications on Pure and Applied Mathematics} \textbf{39} (1986), suppl., S149--S186. DOI: 10.1002/cpa.3160390710.
\bibitem{PLP61pre} A. Pinochet Lobos and C. Pittet, The exact convergence rate in the ergodic theorem of Lubotzky--Phillips--Sarnak, arXiv:1805.05261v2 (2019), Theorem 1.1. Published in expanded form as \emph{The exact convergence rate in the ergodic theorem of Lubotzky--Phillips--Sarnak and a universal lower bound on discrepancies}, \emph{L'Enseignement Math\'ematique} \textbf{67} (2021), nos. 1--2, 63--94. DOI: 10.4171/LEM/1003.
\item[]\textit{Nonhomogeneous Markov chains and observation tails}
\bibitem{Blackwell62} D. Blackwell, Finite non-homogeneous chains, \emph{Annals of Mathematics} (2) \textbf{46} (1945), no. 4, 594--599. DOI: 10.2307/1969199.
\bibitem{Cohn62} H. Cohn, On the tail $\sigma$-algebra of the finite inhomogeneous Markov chains, \emph{Annals of Mathematical Statistics} \textbf{41} (1970), no. 6, 2175--2176. DOI: 10.1214/aoms/1177696725.
\bibitem{GS62} M. Grabchak and I. Sonin, A zero-one law for Markov chains, arXiv:2011.04063 (2020), Theorem 1 and Section 5.1; published in \emph{Stochastics}, DOI: 10.1080/17442508.2021.1980569.
\end{thebibliography}
